% \iffalse meta-comment
%
% banshi.dtx 是排版引擎 banshi 的包头
%
% \fi
%
% \section{排版引擎 \pkg{banshi}}
%
% \changes{v3.2a}{2026/09/23}{照 Word 排版的三个模块（页面设置、段落格式、断行）抽成
%   \pkg{banshi} 包，类在原来的位置加载它}
%
% 学校的范例是 Word 排的。把 Word 的页面设置、段落格式与断行规则照着做出来的那一层，
% 与哈工大哪一份规范都无关，所以单独成包，叫 \pkg{banshi}（版式），与 iota-hit 底下
% 的 Typst 包同名。两边是姊妹实现，键名尽量一一对应，用户在两边的写法可以照着映射。
%
% 眼下它还住在 \hithesis\ 的仓库里：三个模块原样搬进 \file{src/banshi/}，
% 生成 \file{banshi.sty}，类在原先内联它们的位置 \cs{RequirePackage}\marg{banshi}，
% 排出来与搬之前逐字相同。接下来几步照计划做：前缀从 \texttt{hit} 换成
% \texttt{banshi}、学校那几档的元素名改由类注册、消息改英文、\hologo{XeLaTeX} 那条
% 退路也收进来；接口稳了再拆成独立仓库。
%
% ^^A ======================================================================
% ^^A Module package: banshi.sty 的包头
% ^^A ======================================================================
%    \begin{macrocode}
%<*package>
\NeedsTeXFormat{LaTeX2e}[2022/06/01]
\ProvidesExplPackage{banshi}{2026/09/23}{0.1.0}
  {Word-compatible page setup, paragraph formats and line breaking}
%</package>
%    \end{macrocode}
%
% \Finale
